\section{模型假设及适用范围}

题设参数按确定值处理。无人机沿两节点水平直线飞行，巡航高度由DEM净空决定；不考虑风场、临时禁飞区与地面道路通行。不可拆货箱是运输的最小单位，一箱只能在原始需求表指定的服务区交付一次。第一问每架次仅服务一个服务区；后两问允许同架次多点交付，并按交付后的剩余箱重更新能耗。

准备和装载从架次开始时占用实体机及电池；返航结束后实体机可开始下一次准备，电池须先充满。附件未规定共享工位数量，因此不额外施加工位互斥。中继在给定悬停点保持不动。每次任务由O01出发并返回，服务前完成准备和建链；悬停和通信附加功率依附件计入能源消耗。中继任务结束与下一次准备之间保留给定的300~s周转时间。

全部80箱以完整交接结束时刻计实际送达。医疗物资须满足附件期望送达时刻，首批保障货箱须满足首批截止；兼属两类时同时满足。64箱非医疗物资的期望时刻用于统计迟到和配送及时性，不作为不可行判据，其中15箱首批饮用水还受首批截止约束。题设DEM按地形表面使用；其标称高程、节点高程与通信设备高度视为同一垂直基准。若存在定位误差或未表示的新障碍，需重新核算航线和通信裕量。

关于目标权衡，本文先检验全部目的地、医疗及首批时限、安全、资源及通信硬约束，再按优先加权迟到、全部设备最后返航时间和优先加权交付时刻逐级比较，同时报告任务耗电和架次数。有限候选搜索所得“相对最优”是经回代验证的可行改进，不代表全空间全局最优；标称可行也不等同现场天气下的稳健可行。
